


Ce document constitue le rapport de mon stage de Projet de Fin d'Année (PFA), d'une durée de 2 mois. Cette expérience représente une étape charnière de mon cursus académique, offrant une transition pratique vers le monde professionnel par l'application des connaissances acquises en salle de cours à des défis technologiques réels et innovants.

Mon rôle au sein de l'entreprise a été axé sur la participation à un projet ambitieux de développement logiciel nommé "AI-BioProfile". Ce projet visait à créer un système intelligent capable d'automatiser le traitement et la structuration des CV des candidats, répondant ainsi à un besoin critique d'optimisation des processus de recrutement et de présentation commerciale.

Plus spécifiquement, le projet consistait en la construction d'un pipeline d'intelligence artificielle robuste, utilisant des technologies de pointe dans les domaines du Traitement du Langage Naturel (NLP), de la Vision par Ordinateur (OCR) et des Grands Modèles de Langage (LLM) exécutés localement pour garantir la souveraineté des données. Ce travail a impliqué la conception et la mise en place d'une architecture permettant d'extraire intelligemment les informations non structurées des CV, et de générer dynamiquement des dossiers de compétences standardisés (BioProfiles) sous format PowerPoint, prêts à être présentés aux clients.

Durant mon stage, j'ai été amené(e) à :
\begin{itemize}
    \item Analyser les besoins spécifiques des Ingénieurs d'Affaires pour développer une solution logicielle sur mesure.
    \item Concevoir et développer une architecture de données intégrant l'extraction multimodale (texte et image) via OCR et l'inférence par Intelligence Artificielle locale.
    \item Implémenter la génération dynamique et automatisée de documents commerciaux interactifs (PowerPoint).
    \item Intégrer une validation humaine (Human-in-the-loop) permettant aux utilisateurs de vérifier et corriger les données extraites par l'IA.
    \item Tester et optimiser les performances des modèles et du pipeline pour garantir une précision, une fiabilité et une confidentialité maximales des données personnelles.
\end{itemize}

Ce rapport détaille chacune de ces étapes, offrant une vue complète du processus de développement, des défis techniques rencontrés (notamment l'hétérogénéité des formats de CV), des solutions implémentées, ainsi que des résultats obtenus et de leur impact sur la productivité des équipes. L'expérience acquise durant cette période a été extrêmement enrichissante, non seulement en termes d'acquisition de compétences techniques avancées en IA, mais aussi pour mon développement professionnel et personnel.
